\title{Group Sequence Policy Optimization}

\begin{abstract}
We present Group Sequence Policy Optimization (GSPO), a reinforcement learning algorithm that is both stable and effective for training large language models. Distinct from earlier methods that rely on token-level importance ratios, GSPO operates by establishing the importance ratio at the sequence likelihood level and executes clipping, reward assignment, and optimization over entire sequences. Our results show that GSPO not only surpasses the GRPO algorithm in terms of training speed and overall performance, but also significantly enhances the stability of Mixture-of-Experts (MoE) RL training and holds promise for streamlining RL infrastructure. These advantages have played a crucial role in achieving the outstanding improvements observed in the latest Qwen3 models.
\end{abstract}

\section{Introduction}
\label{sec:intro}

Reinforcement learning (RL) has become an essential framework for the advancement of large-scale language models \citep{o1,dpsk-r1,qwq32b,qwen3}. By leveraging extensive RL, language models acquire the proficiency to solve highly complex tasks—including those in advanced mathematics and programming—by engaging in deeper, more intricate reasoning.

To effectively scale RL by increasing computational resources, it is critical to preserve stable and reliable training dynamics. Yet, current leading RL techniques, such as GRPO \citep{grpo}, are plagued by significant instability when applied to enormous language models; this often leads to disastrous and irreversible degradation of the model’s performance \citep{qwen3, minimax-m1}. Such instability poses a major obstacle to enhancing the abilities of language models through ongoing RL-based training.

This study demonstrates that the core cause of instability in GRPO is the improper use and subsequent invalidation of importance sampling weights during the algorithm's operation. This flaw generates substantial training noise with high variance, which escalates as response sequences grow longer and becomes even more pronounced due to the clipping procedure, ultimately leading to the collapse of the model.

To overcome these fundamental challenges, we introduce \textbf{Group Sequence Policy Optimization (GSPO)}, a novel reinforcement learning framework for optimizing large language models. GSPO's central contribution is its rigorously justified approach to defining importance ratios, grounded in sequence likelihood as established in \citep{click}, and thus consistent with the foundational principles of importance sampling. Furthermore, GSPO calculates normalized rewards by treating the comparative advantages among multiple responses to a single query, guaranteeing coherent optimization at the sequence reward level.

Through our empirical studies, we show that GSPO markedly surpasses GRPO in terms of training robustness, speed, and overall effectiveness. Notably, GSPO intrinsically addresses the instability issues commonly encountered during reinforcement learning of large Mixture-of-Experts (MoE) models, removing the dependence on intricate stabilization techniques and offering a path toward more streamlined RL pipelines. These advantages have directly facilitated substantial performance gains in the development of the latest Qwen3 models. We anticipate that GSPO will serve as a reliable and scalable algorithmic cornerstone for further progress in large-scale RL training for language models.

\section{Preliminaries}

\paragraph{Notation}
Throughout this work, we characterize an autoregressive language model with parameters $\theta$ as a policy $\pi_\theta$.
The variable $x$ signifies a query, while $\mathcal{D}$ refers to the collection of all queries. For any response $y$ corresponding to a particular query $x$, the probability assigned by the policy $\pi_\theta$ is represented as $\pi_\theta (y | x)=\prod_{t=1}^{|y|} \pi_\theta (y_t | x, y_{<t} )$, with $|y|$ denoting the token count in $y$. Each query-response pair $(x, y)$ is evaluated by a verifier function $r$, which generates a reward $r(x, y) \in [ 0, 1 ]$.

\paragraph{Proximal Policy Optimization (PPO)} Using samples generated from the old policy $\pi_{\theta_\text{old}}$, PPO \citep{ppo} constrains the policy update within a proximal region of the old policy through the clipping mechanism. Specifically, PPO employs the following objective for policy optimization (we omit the KL regularization term hereinafter for brevity, as it is not the focus of this paper):

\begin{align}
\mathcal{J}_\text{PPO}(\theta) = \mathbb{E}_{ x \sim \mathcal{D},\, y \sim \pi_{\theta_\text{old}}( \cdot | x) }
\left[ \frac{1}{|y|} \sum_{t=1}^{|y|} 
\min \left( w_{t}(\theta) \widehat{A}_{t},  \, \mathrm{clip} \left( w_{t}(\theta), 1 - {\varepsilon}, 1 + {\varepsilon}\right) \widehat{A}_{t} \right)
\right],
\end{align}

Here, the importance ratio for token $y_t$ is given by $
w_{t}(\theta) = \frac{ \pi_{\theta} (y_{t} | x, y_{<t}) }{ \pi_{\theta_\text{old}} (y_{t} | x,y_{<t})} $, where $\widehat{A}_{t}$—the advantage of $y_t$—is determined using a separate value model, and $\varepsilon$ denotes the range for clipping importance ratios.

In practice, PPO's primary obstacle is its significant dependence on the value model. The value model is often comparable in scale to the policy model, resulting in substantial demands for memory and computation. The effectiveness of PPO is tightly linked to the precision of value estimates, which is inherently difficult to attain. The task becomes even more arduous when aiming for robust value estimates that are applicable to extended responses or tasks of greater complexity.

\paragraph{Group Relative Policy Optimization (GRPO)} GRPO \citep{grpo} bypasses the need for the value model by computing the relative advantage of each response within a group of responses to the same query. Specifically, GRPO optimizes the following objective:

\begin{align}
\label{equ:grpo}
\mathcal{J}_\text{GRPO}(\theta) = \mathbb{E}_{ x \sim \mathcal{D},\, \{y_i\}_{i=1}^G \sim \pi_{\theta_\text{old}}( \cdot | x) }
\left[ \frac{1}{G} \sum_{i=1}^{G} \frac{1}{|y_i|} \sum_{t=1}^{|y_i|} 
\min \left( w_{i,t}(\theta) \widehat{A}_{i,t},  \, \mathrm{clip} \left( w_{i,t}(\theta), 1 - {\varepsilon}, 1 + {\varepsilon}\right) \widehat{A}_{i,t} \right)
\right],
\end{align}

where $G$ is the number of generated responses to each query $x$ (i.e., the group size), and the importance ratio $w_{i,t}(\theta)$ and advantage $\widehat{A}_{i,t}$ of token $y_{i,t}$ are:

\begin{align}
    w_{i,t}(\theta)=\frac{ \pi_{\theta} (y_{i,t} | x, y_{i,<t}) }{ \pi_{\theta_\text{old}} (y_{i,t} | x,y_{i,<t})},\quad\ 
    \widehat{A}_{i,t} = \widehat{A}_{i} = \frac{r(x, y_i) - \mathrm{mean} \left( \{ r(x, y_i) \}_{i=1}^G \right) }{ \mathrm{std} \left( \{ r(x, y_i) \}_{i=1}^G \right) },
\end{align}

respectively, where all the tokens in $y_i$ share the same advantage as $\widehat{A}_{i}$.

\section{Motivation}
\label{sec:motivation}

As model architectures expand in size and complexity—such as through increased sparsity in Mixture-of-Experts approaches—and with longer generated responses, larger rollout batch sizes are required in reinforcement learning to effectively leverage hardware resources. In order to enhance sample efficiency, it is customary to divide these extensive rollout batches into several mini-batches for performing gradient updates. Such a strategy inherently results in an off-policy scenario: the generated outputs $y$ originate from a previous policy $\pi_{\theta_\text{old}}$, distinct from the current policy $\pi_\theta$ undergoing optimization. Consequently, this situation highlights the importance of the clipping mechanisms found in both PPO and GRPO, which serve to exclude samples that are excessively "off-policy" from contributing to gradient calculations.

While mechanisms like clipping aim to manage this off-policy discrepancy, we identify a more fundamental issue in GRPO: \textit{its objective is ill-posed}. This problem becomes particularly acute when training large models on long-response tasks, leading to catastrophic model collapse. The ill-posed nature of the GRPO objective stems from a misapplication of importance sampling weights. The principle of importance sampling is to estimate the expectation of a function $f$ under a target distribution $\pi_\text{tar}$ by re-weighting samples drawn from a behavior distribution $\pi_\text{beh}$:

\begin{align}
\mathbb{E}_{z \sim \pi_\text{tar}} \left[ f(z) \right] = \mathbb{E}_{z \sim \pi_\text{beh}} \left[ \frac{ \pi_\text{tar} (z) }{ \pi_\text{beh} (z)} \, f(z) \right].
\end{align}

Fundamentally, effective correction for differences between distributions using the importance weight $\frac{ \pi_\text{tar} (z) }{ \pi_\text{beh} (z)}$ depends on computing an average across a large number of samples ($N \gg 1$) drawn from the behavior policy $\pi_\text{beh}$. 

On the other hand, GRPO computes the importance weight $\frac{ \pi_{\theta} (y_{i,t} | x, y_{i,<t}) }{ \pi_{\theta_\text{old}} (y_{i,t} | x,y_{i,<t})}$ at every token $t$, using only a single token sample $y_{i,t}$ per next-token distribution $\pi_{\theta_\text{old}}( \cdot | x, y_{i, <t})$. This strategy does not effectively correct for the underlying distribution shift. Instead, it increases variance in the gradient estimates, which accumulates as the sequence grows longer, especially because of the use of clipping. Our experiments show that such variance can trigger an irreversible collapse in model performance. If collapse happens, further training cannot restore the model, even when rolling back to an earlier checkpoint, tweaking hyperparameters like the clipping thresholds, increasing the sequence length for generation, or altering the RL query structure.

This empirical outcome exposes a key flaw within the GRPO framework. The breakdown of importance weighting at the token level underscores a fundamental insight: \textbf{optimization granularity should align with the granularity of reward assignment}. Since rewards are associated with the full sequence, attempting off-policy correction at individual token positions is ill-suited. Consequently, this rationale leads us to abandon token-level objectives in favor of computing importance weights and conducting optimization at the \textit{sequence level}.

\section{Algorithm}

\subsection{GSPO: Group Sequence Policy Optimization}

Although the use of the token-level importance weight $\frac{ \pi_{\theta} (y_{i,t} | x, y_{i,<t}) }{ \pi_{\theta_\text{old}} (y_{i,t} | x,y_{i,<t})}$ poses challenges in GRPO, we note that the \textit{sequence-level} importance weight $\frac{ \pi_{\theta} (y | x) }{ \pi_{\theta_\text{old}} (y | x)}$ carries a well-defined theoretical interpretation in language generation tasks. Specifically, it quantifies the extent to which a response $y$—sampled from $\pi_{\theta_\text{old}} (\cdot | x)$—diverges from the distribution $\pi_{\theta} (\cdot | x)$. This measure not only fits naturally with sequence-level rewards but also acts as a relevant metric for the clipping process.

Based on this straightforward observation, we propose the \textbf{Group Sequence Policy Optimization (GSPO)} algorithm. GSPO employs the following sequence-level optimization objective:

\begin{align}
\mathcal{J}_\text{GSPO} (\theta) =
\mathbb{E}_{ x \sim \mathcal{D},\, \{y_i\}_{i=1}^G \sim \pi_{\theta_\text{old}}( \cdot | x) }
\left[ 
\frac{1}{G} \sum_{i=1}^{G}
\min \left( s_{i}(\theta)  \widehat{A}_{i},  \, \mathrm{clip} \left( s_{i}(\theta), 1 - {\varepsilon}, 1 + {\varepsilon} \right) \widehat{A}_{i} \right) 
\right],
\label{equ:gspo}
\end{align}

where we adopt the group-based advantage estimation:

\begin{align}
\widehat{A}_{i} = \frac{r(x, y_i) - \mathrm{mean} \left( \{ r(x, y_i) \}_{i=1}^G \right) }{ \mathrm{std} \left( \{ r(x, y_i) \}_{i=1}^G \right) },
\end{align}

and define the importance ratio $s_{i}(\theta)$ based on sequence likelihood \citep{click}:

\begin{align}
s_{i}(\theta) = \left( \frac{ \pi_{\theta} (y_i | x) }{ \pi_{\theta_\text{old}} (y_i | x)} \right)^{\frac{1}{|y_i|}}
=
\exp \left( \frac{1}{|y_i|} \sum_{t=1}^{|y_i|} \log \frac{ \pi_{\theta} (y_{i,t} | x, y_{i,<t}) }{ \pi_{\theta_\text{old}} (y_{i,t} | x,y_{i,<t})} \right).
\end{align}

As a result, GSPO utilizes clipping at the response level rather than at the token level, ensuring that highly ``off-policy'' samples are removed from the gradient calculation, thereby aligning both the rewarding and optimization processes with sequence-level criteria. Additionally, length normalization is incorporated into $s_{i}(\theta)$ to stabilize variance and constrain $s_{i}(\theta)$ within a consistent numerical interval. Without this normalization, minor changes in token likelihood can cause substantial swings in the sequence-level importance ratio, necessitating different clipping thresholds for responses of varying lengths. Furthermore, we observe that the clipping intervals employed in GSPO typically have a markedly different scale compared to prior approaches such as GRPO, owing to the unique manner in which importance ratios are defined.

\subsection{Gradient Analysis}
\label{subsec:gradient}

We can derive the gradient of the GSPO objective as follows (clipping is omitted for brevity):

\begin{align}
\nabla_{\theta} \mathcal{J}_\text{GSPO} (\theta)
=&\ 
\nabla_{\theta} \mathbb{E}_{ x \sim \mathcal{D},\, \{y_i\}_{i=1}^G \sim \pi_{\theta_\text{old}}( \cdot | x) }
\left[ 
\frac{1}{G} \sum_{i=1}^{G} s_{i}(\theta) \widehat{A}_{i}
\right] \\
= &\ 
\mathbb{E}_{ x \sim \mathcal{D},\, \{y_i\}_{i=1}^G \sim \pi_{\theta_\text{old}}( \cdot | x) }
\left[ 
\frac{1}{G} \sum_{i=1}^{G}
s_{i}(\theta) \widehat{A}_{i}
\cdot \nabla_{\theta} \log s_{i}(\theta)
\right] \\
=&\ 
\mathbb{E}_{ x \sim \mathcal{D},\, \{y_i\}_{i=1}^G \sim \pi_{\theta_\text{old}}( \cdot | x) }
\left[ 
\frac{1}{G} \sum_{i=1}^{G} 
\left( \frac{ \pi_{\theta} (y_i | x) }{ \pi_{\theta_\text{old}} (y_i | x)} \right)^{\frac{1}{|y_i|}} \widehat{A}_{i} 
\cdot \frac{1}{|y_i|} \sum_{t=1}^{|y_i|} \nabla_{\theta} \log \pi_{\theta} (y_{i,t} | x, y_{i,<t}) 
\right].
\label{equ:gspo_grad}
\end{align}

For comparison, the gradient of the GRPO objective is as follows (note that $\widehat{A}_{i,t} = \widehat{A}_{i}$):

\begin{align}
\nabla_{\theta} \mathcal{J}_\text{GRPO} (\theta) 
=&\ 
\nabla_{\theta} \mathbb{E}_{ x \sim \mathcal{D},\, \{y_i\}_{i=1}^G \sim \pi_{\theta_\text{old}}( \cdot | x) }
\left[ \frac{1}{G} \sum_{i=1}^{G} \frac{1}{|y_i|} \sum_{t=1}^{|y_i|} 
w_{i,t}(\theta) \widehat{A}_{i,t}
\right] \\
=&\
\mathbb{E}_{ x \sim \mathcal{D},\, \{y_i\}_{i=1}^G \sim \pi_{\theta_\text{old}}( \cdot | x) }
\left[ \frac{1}{G} \sum_{i=1}^{G} \widehat{A}_{i} 
\cdot \frac{1}{|y_i|} \sum_{t=1}^{|y_i|} 
\frac{ \pi_{\theta} (y_{i,t} | x, y_{i,<t}) }{ \pi_{\theta_\text{old}} (y_{i,t} | x,y_{i,<t})} 
\nabla_{\theta} \log \pi_{\theta} (y_{i,t} | x, y_{i,<t})  
\right].
\end{align}

The key difference between GSPO and GRPO concerns \textit{the method used to assign weights to the gradients of token log likelihoods}. GRPO employs token-specific ``importance weights'' given by $\frac{ \pi_{\theta} (y_{i,t} | x, y_{i,<t}) }{ \pi_{\theta_\text{old}} (y_{i,t} | x,y_{i,<t})}$, which are not uniformly distributed. These varying weights may range over $(0, 1+\varepsilon]$ when $\widehat{A}_{i} >0$, or $[1-\varepsilon, +\infty)$ if $\widehat{A}_{i} < 0$, and thus their effects are substantial, potentially accumulating during training and generating erratic results. Conversely, GSPO applies uniform weighting to all tokens within a response, removing the instability introduced by the unequal weighting scheme in GRPO.

\subsection{GSPO-token: A Token-level Objective Variant}

In scenarios like multi-turn RL, we may desire a finer-grained advantage adjustment than the sequence level. To this end, we introduce a token-level objective variant of GSPO, namely \textbf{GSPO-token}, to allow token-wise advantage customization:

\begin{align}
\mathcal{J}_\text{GSPO-token}(\theta)
=
\mathbb{E}_{ x \sim \mathcal{D},\, \{y_i\}_{i=1}^G \sim \pi_{\theta_\text{old}}( \cdot | x) }
\left[ 
\frac{1}{G} \sum_{i=1}^{G} \frac{1}{|y_i|} \sum_{t=1}^{|y_i|} 
\min \left( s_{i,t}(\theta) \widehat{A}_{i,t},  \, \mathrm{clip} \left( s_{i,t}(\theta), 1 - {\varepsilon}, 1 + {\varepsilon} \right) \widehat{A}_{i,t} \right) 
\right],
\label{equ:gspo-token}
\end{align}

where

\begin{align}
s_{i,t}(\theta) 
= \mathrm{sg} \left[ s_{i}(\theta) \right]  \cdot \frac{ \pi_{\theta} (y_{i,t} | x, y_{i,<t}) }{ \mathrm{sg} \left[ \pi_{\theta} (y_{i,t} | x, y_{i,<t}) \right] },
\end{align}

and $\mathrm{sg}[\cdot]$ denotes only taking the numerical value but stopping the gradient, corresponding to the \texttt{detach} operation in PyTorch. The gradient of GSPO-token can be derived as:

\begin{align}
\nabla_{\theta} \mathcal{J}_\text{GSPO-token}(\theta)
=&\ 
\nabla_{\theta} \mathbb{E}_{ x \sim \mathcal{D},\, \{y_i\}_{i=1}^G \sim \pi_{\theta_\text{old}}( \cdot | x) }
\left[ 
\frac{1}{G} \sum_{i=1}^{G}
\frac{1}{|y_i|} \sum_{t=1}^{|y_i|} 
s_{i,t}(\theta) \widehat{A}_{i,t}
\right] \\
=&\ 
\mathbb{E}_{ x \sim \mathcal{D},\, \{y_i\}_{i=1}^G \sim \pi_{\theta_\text{old}}( \cdot | x) }
\left[ 
\frac{1}{G} \sum_{i=1}^{G} s_i (\theta) \cdot \frac{1}{|y_i|} \sum_{t=1}^{|y_i|} 
\widehat{A}_{i,t} \frac{ \nabla_{\theta} \pi_{\theta} (y_{i,t} | x, y_{i,<t}) }{ \pi_{\theta} (y_{i,t} | x, y_{i,<t}) }
\right] \\
=&\ 
\mathbb{E}_{ x \sim \mathcal{D},\, \{y_i\}_{i=1}^G \sim \pi_{\theta_\text{old}}( \cdot | x) }
\left[ 
\frac{1}{G} \sum_{i=1}^{G}  \left( \frac{ \pi_{\theta} (y_i | x) }{ \pi_{\theta_\text{old}} (y_i | x)} \right)^{\frac{1}{|y_i|}}
\cdot \frac{1}{|y_i|} \sum_{t=1}^{|y_i|}
\widehat{A}_{i,t} \nabla_{\theta} \log \pi_{\theta} (y_{i,t} | x, y_{i,<t})  
\right].
\label{equ:gspo-token_grad}
\end{align}

It should be observed that the expression $\frac{ \pi_{\theta} (y_{i,t} | x, y_{i,<t}) }{ \mathrm{sg} \left[ \pi_{\theta} (y_{i,t} | x, y_{i,<t}) \right] }$ always evaluates to 1, which implies that $s_{i,t}(\theta)$ is, in fact, equivalent in value to $s_{i}(\theta)$. By inspecting Equation~\eqref{equ:gspo} alongside Equation~\eqref{equ:gspo-token}, as well as Equation~\eqref{equ:gspo_grad} compared with Equation~\eqref{equ:gspo-token_grad}, it becomes clear that GSPO-token and GSPO yield identical numerical outcomes in terms of their optimization targets, clipping rules, and analytic gradients—given the condition where each token’s advantage in the sequence $y_i$ is set to an identical value (i.e., $\widehat{A}_{i,t} = \widehat{A}_{i}$). In contrast, GSPO-token provides increased adaptability by allowing distinct advantage assignments for each token.

\section{Experiments and Discussion}

\subsection{Empirical Results}

We conduct experiments using a cold-start model initialized from Qwen3-30B-A3B-Base, presenting both training reward trajectories and evaluation metrics on several benchmarks: AIME'24 (mean Pass@1 across 32 samples), LiveCodeBench (202410-202502, mean Pass@1 over 8 samples), and CodeForces (Elo Rating). Throughout reinforcement learning training, each rollout batch is subdivided into four mini-batches for the purpose of gradient computation. For GSPO, the left and right clipping thresholds in Equation~\eqref{equ:gspo} are configured at 3e-4 and 4e-4, respectively. In our baseline comparison with GRPO, we set the left and right clipping bounds in Equation~\eqref{equ:grpo} at 0.2 and 0.27, carefully selected for equitable evaluation. It is important to mention that successful RL convergence with GRPO in the Mixture of Experts framework requires the Routing Replay strategy, as detailed in \S~\ref{subsec:routing_replay}; \textit{however, GSPO eliminates the necessity for this additional training procedure}.

Figure~\ref{fig:results} illustrates that training with GSPO remains stable throughout the process. Our results indicate that \textbf{GSPO consistently enhances performance by allocating greater training compute, periodically refreshing the query set, and lengthening the generated sequences}. Additionally, GSPO outperforms GRPO in terms of training efficiency, yielding superior training accuracy and benchmark results when given identical training compute and number of queries. Lastly, we have implemented GSPO in the RL training of the most recent Qwen3 models, providing compelling evidence of GSPO's effectiveness in harnessing the advantages of RL scaling for large language models.

\subsection{Curious Observation on Clipping Fractions}

GSPO differs fundamentally from GRPO in that it applies clipping to entire responses rather than to each token separately. As illustrated in Figure~\ref{fig:clipping}, the proportion of clipped tokens in GSPO versus GRPO differs by approximately two orders of magnitude, and this relationship remains consistent regardless of the chosen clipping ranges. Interestingly, although GSPO clips many more tokens—resulting in a reduced set of tokens utilized for training or gradient estimation—it still outperforms GRPO in terms of training efficiency. This unexpected result, in which extensive token clipping correlates with enhanced training efficiency, implies that token-level gradient estimates in GRPO are prone to noise and fail to effectively exploit samples. Conversely, GSPO's sequence-level clipping yields a more stable and valuable learning signal.

\subsection{Benefit of GSPO for MoE Training}
\label{subsec:routing_replay}

\paragraph{Background}
In contrast to reinforcement learning (RL) applied to conventional dense models, the sparse activation pattern characteristic of MoE architectures presents distinct challenges with respect to training stability. Specifically, our observations with the GRPO algorithm indicate that \textit{expert-activation volatility} within MoE models can significantly disrupt the convergence of RL training. More concretely, following one or more gradient steps, the set of experts engaged to generate an identical output can alter substantially. To illustrate, experiments with the 48-layer Qwen3-30B-A3B-Base configuration reveal that, on average, approximately 10\% of the experts selected for a given rollout differ post-update under the updated policy $\pi_{\theta}$ relative to those activated under the preceding policy $\pi_{\theta_\text{old}}$.
This tendency, which escalates with increasing MoE model depth, results in pronounced fluctuations in the token-level importance ratios $w_{i,t}(\theta) = \frac{ \pi_{\theta} (y_{i,t} | x, y_{i,<t}) }{ \pi_{\theta_\text{old}} (y_{i,t} | x,y_{i,<t})}$, ultimately compromising their reliability as detailed in \S~\ref{sec:motivation} and~\ref{subsec:gradient}. The consequent instability thus poses a major obstacle to the successful convergence of RL-based optimization.

\paragraph{Our Previous Approach} In earlier work, we adopted the \textbf{Routing Replay} methodology to address this problem. The approach entails recording the experts activated under $\pi_{\theta_\text{old}}$ and then reusing these activation patterns within $\pi_{\theta}$ during computation of the importance ratios $w_{i,t}(\theta) = \frac{ \pi_{\theta} (y_{i,t} | x, y_{i,<t}) }{ \pi_{\theta_\text{old}} (y_{i,t} | x,y_{i,<t})}$. By enforcing identical expert activation for both $\pi_{\theta} (y_{i,t} | x, y_{i,<t})$ and $\pi_{\theta_\text{old}} (y_{i,t} | x,y_{i,<t})$ for each token $y_{i,t}$, this mechanism helps maintain the reliability of token-level importance ratios. Consequently, optimization remains focused on a consistent activated subnet across successive gradient steps. As illustrated in Figure~\ref{fig:routing_replay}, Routing Replay is pivotal for achieving stable GRPO training dynamics in MoE architectures.

\paragraph{Benefit of GSPO} While Routing Replay allows MoE models trained with GRPO to achieve stable convergence, its reliance on reusing routing patterns results in increased demands on memory and inter-node communication, and it may restrict the available capacity of the MoE architecture. In comparison, as depicted in Figure~\ref{fig:results}, GSPO bypasses the necessity for Routing Replay by directly calculating the importance ratios $s_i(\theta)$ in a standard manner, achieving stable optimization and convergence. The central observation is that GSPO’s performance depends exclusively on sequence-level likelihood (i.e., $\pi_{\theta} (y_i | x)$) rather than on token-level likelihoods (i.e., $\pi_{\theta} (y_{i,t} | x, y_{i,<t})$), ensuring robustness. Because MoE models consistently preserve strong language modeling, the likelihood over entire sequences remains relatively stable. Thus, GSPO inherently solves the expert-activation instability in MoE frameworks, eliminating the necessity for intricate approaches such as Routing Replay. This leads to a more straightforward and reliable training workflow and enables the model to fully utilize its capacity, free from artificial limitations.

\subsection{Benefit of GSPO for RL Infrastructure}

Owing to the fact that training engines such as Megatron and inference engines like SGLang and vLLM exhibit differences in numerical precision, practitioners ordinarily rely on the training engine to recalculate the likelihoods of generated responses under the old policy $\pi_{\theta_\text{old}}$. In contrast, GSPO employs only the sequence-level likelihoods in its optimization process, as opposed to token-level ones, and thus, in theory, is significantly more robust to precision mismatches. As a result, GSPO enables practitioners to leverage the likelihoods outputted directly from the inference engine during optimization, thereby circumventing the necessity of recalculation using the training engine. This advantage becomes particularly pronounced in cases such as partial rollout, multi-turn RL, and frameworks where training and inference operations are decoupled.

\section{Conclusion}

Group Sequence Policy Optimization (GSPO) is introduced as an innovative reinforcement learning method tailored for training large language models. GSPO leverages importance sampling by utilizing sequence likelihood to construct importance ratios, implementing sequence-level clipping, rewarding, and optimization. Compared to GRPO, GSPO offers significant improvements in training stability, efficiency, and overall performance. It proves especially effective in large-scale reinforcement learning for MoE models, contributing substantially to the progress seen in recent Qwen3 models. Positioned as a scalable algorithmic foundation, GSPO paves the way for further scaling in reinforcement learning, with anticipation for transformative developments in the field of artificial intelligence.

\end{document}